% --- 1. INFORMAZIONI GENERALI ---
\section{Informazioni Generali}
\begin{itemize}
	\item \textbf{Luogo/Piattaforma:} {Server Discord\textsubscript{G}}
	\item \textbf{Data:} 2026-04-21
	\item \textbf{Orario di inizio:} 13:30
	\item \textbf{Orario di fine:} 14:45
	\item \textbf{Responsabile\textsubscript{G}:} {Francesco Marcon}
	\item \textbf{Scriba:} {Sofia De Blasi}
\end{itemize}

\vspace{0.5cm}
\textbf{Partecipanti presenti:}
\begin{table}[ht]
	\centering
	\renewcommand{\arraystretch}{1.2}
	\begin{tabularx}{\textwidth}{X l}
		\toprule
		\textbf{Nome e Cognome} & \textbf{Matricola} \\
		\midrule
		Sofia De Blasi & 2111014 \\
		Roberto Mariano Doroftei & 2111031 \\
		Filippo Idiometri & 2035617 \\
		Francesco Marcon & 2110990 \\
		Armando Moda Scarati & 2082864 \\
		Anton Ricardo Rupi & 2054304 \\
		\bottomrule
	\end{tabularx}
\end{table}

\vspace{0.5cm}
\textbf{Partecipanti assenti:}
\begin{table}[ht]
	\centering
	\renewcommand{\arraystretch}{1.2}
	\begin{tabularx}{\textwidth}{X l}
		\toprule
		\textbf{Nome e Cognome} & \textbf{Matricola} \\
		\midrule
		Nessuno & - \\
		\bottomrule
	\end{tabularx}
\end{table}

% --- 2. ORDINE DEL GIORNO ---
\section{Ordine del Giorno}
Gli argomenti previsti per la riunione odierna sono:
\begin{enumerate}
	\item Sprint\textsubscript{G} review e retrospective
	\item Messa in pari sulla stesure dell'AdR e del PdP
	\item Pianificazione dello sprint successivo e assegnazione dei task\textsubscript{G} e ruoli
\end{enumerate}

% --- 3. RESOCONTO E DISCUSSIONE ---
\section{Resoconto della Riunione}
\subsection{Sprint\textsubscript{G} review e retrospective}
Il gruppo si è confrontato sull'andamento dello sprint appena terminato. Sono state segnalate delle criticità su un'aggiunta effettuata al documento di Analisi dei Requisiti, specificando che, l'aggiunta di uno Use Case non previsto, comporta una modifica di tutti gli indici degli Use Case successivi. Questo, durante lo sprint, ha causato un leggero aumento dei tempi previsti, in quanto, il ruolo di Analista, ha dovuto non solo aggiornare tutti gli indici a cascata (dopo l'aggiunta dello Use Case di "Eliminazione di Nota"), ma anche ricreare tutti i diagrammi per aggiornare gli indici di riferimento al loro interno. A questo proposito, è stata assunta la decisione di mantenere una copia di tutti i diagrammi sul software Lucid Spark, così che, in caso debbano essere aggiunti ulteriori Use Case in futuro, anche i diagrammi possano essere facilmente e rapidamente modificati, senza essere ricreati da zero. \\
La stesura del Piano di Progetto ha richiesto più tempo del previsto, 5 ore totali, quando ne erano state previste 2. Questo ha fatto si che il gruppo si interrogasse sul reale dispendio di tempo previsto dal documento del Piano di Progetto, ma ne è emerso che il tempo così prolungato è stato necessario per definire più specificatamente le sezioni del documento, e infine per predisporre un template da utilizzare per la futura stesura delle sezioni ripetitive che ad ogni sprint devono essere compilate.

\subsection{Messa in pari sulla stesure dell'AdR e del PdP}
Durante la riunione il gruppo ha effettuato un 
allineamento sullo stato di avanzamento 
dell’Analisi dei Requisiti\textsubscript{G} (AdR) e del Piano 
di Progetto (PdP), con l’obiettivo di garantire 
una visione condivisa del lavoro svolto e delle 
attività da proseguire nello sprint successivo. 
È emersa la necessità di chiarire la suddivisione 
delle responsabilità, soprattutto per documenti 
estesi come l’AdR, che richiedono un confronto 
collettivo prima della ripartizione delle sezioni. 
Il team ha inoltre discusso 
l’importanza di mantenere aggiornati i 
documenti in modo coerente, definendo convenzioni 
comuni e assicurando che ogni 
membro comprenda lo stato attuale dei contenuti. 
L’allineamento ha permesso di chiarire dubbi, 
uniformare l’approccio redazionale e stabilire le 
priorità per la prosecuzione coordinata dei due documenti.

Oltre alle criticità descritte al paragrafo precedente, la stesura di Adr (Analisi dei Requisiti) e PdP (Piano di Progetto) è terminata con tutte le aggiunte previste dagli scorsi task. La divisione del lavoro tra più membri del gruppo, tramite sezioni del documento, ha permesso una stesura efficace di entrambi i documenti.

\subsection{Pianificazione dello sprint successivo e assegnazione dei task e ruoli}
Lo sprint è stato definito con durata dal 21 al 28 aprile 2026. I ruoli assegnati per questo sprint sono: Responsabile (Francesco Marcon), Amministratore (Sofia De Blasi), Verificatore (Roberto M. Doroftei), Analisti (Filippo Idiometri, Armando Moda Scarati, Anton R. Rupi).\\
Come nello sprint precedente, i ruoli di Programmatore e Progettista non sono ancora necessari, si è infatti deciso di dedicare 3 membri del gruppo al ruolo di Analista, più specificatamente Armando e Ricardo continueranno la stesura dell'Analisi dei Requisiti, in particolare della terza macroarea, con anche gli Use Cases 15 e 19 (opzionali) per i quali sarà avviata una comunicazione via mail, con la proponente Zucchetti, per richiedere ulteriori chiarimenti; mentre Filippo continua la stesura del Piano di Progetto, inserendo il nuovo sprint.
 

% --- 4. DECISIONI PRESE ---
\section{Decisioni Prese}

\renewcommand{\arraystretch}{1.3}
\vspace{-0.3cm}
\noindent
\begin{longtable}{c p{12cm}}
	\toprule
	\textbf{Codice} & \textbf{Decisione} \\
	\midrule
	\endfirsthead

	\toprule
	\textbf{Codice} & \textbf{Decisione} \\
	\midrule
	\endhead

	\midrule
	\multicolumn{2}{r}{\textit{Continua nella pagina successiva}} \\
	\midrule
	\endfoot

	\bottomrule
	\endlastfoot

	\textbf{VI-10.1} & Mantenere una copia dei diagrammi dei Casi d'Uso su LucidSpark per facilitarne la modifica futura. \\[4pt]
	\textbf{VI-10.2} & Definire una modalità di gestione delle ore nel PdP basata su proporzioni (minuti → ore → costo per ruolo). \\[4pt]
	\textbf{VI-10.3} & Iniziare l’analisi della qualità di prodotto e di processo (argomenti T7 e T8) in preparazione al Piano di Qualifica. \\[4pt]
	\textbf{VI-10.4} & Uniformare la nomenclatura dei casi d’uso (formato “UC X: testo”). \\[4pt]
	\textbf{VI-10.5} & Pianificare l’invio di una mail al professor Cardin per organizzare una riunione successiva al 28/04. \\[4pt]

	

	

\end{longtable}

% --- 5. AZIONI (TODO) ---
\section{Azioni da Svolgere (To-Do)}

\renewcommand{\arraystretch}{1.3}
\begin{longtable}{c c p{4.5cm} p{3cm} l}
    \toprule
    \textbf{Codice} & \textbf{Rif. Decisione} & \textbf{Descrizione Task} & \textbf{Assegnatario} & \textbf{Scadenza} \\
    \midrule
    \endfirsthead

    \toprule
    \textbf{Codice} & \textbf{Rif. Decisione} & \textbf{Descrizione Task} & \textbf{Assegnatario} & \textbf{Scadenza} \\
    \midrule
    \endhead

    \midrule
    \multicolumn{5}{r}{\textit{Continua nella pagina successiva}} \\
    \midrule
    \endfoot

    \bottomrule
    \endlastfoot

	% Issue #121 - Aggiornamento Sito
    \ghissue{121} & - & Aggiornare il sito web con i documenti più recenti. & Sofia De Blasi & 2026-04-28 \\[4pt]
	
	% Issue #122 - AdR Casi d'Uso Macroarea 3 + UC15, UC19
    \ghissue{122} & - & Stesura Casi d'Uso UC15, UC19, e tutta la macroarea 3. & Armando Moda Scarati, Anton Ricardo Rupi & 2026-04-28 \\[4pt]

	% Issue #123 - PdP stesura sprint 2 e approvazione per ingresso in baseline
    \ghissue{123} & - & Stesura dello sprint 2 nel PdP, approvazione per ingresso in baseline. & Filippo Idriometri & 2026-04-28 \\[4pt]

    % Issue #124 - Verbale interno
    \ghissue{124} & - & Stesura verbale interno della riunione odierna. & Francesco Marcon & 2026-04-28 \\[4pt]

	% Issue #125 - Verbale interno
    \ghissue{125} & - & Stesura verbale esterno per comunicazione con Zucchetti, e invio mail per richiesta firma & Sofia De Blasi & 2026-04-28 \\[4pt]

	% Issue #126 - Invio mail a Cardin
    \ghissue{126} & VI-10.5 & Invio mail per richiesta ricevimento a prof Cardin & Francesco Marcon & 2026-04-28 \\[4pt]

	% Issue #127 - Riordino diagrammi su lucid spark
    \ghissue{127} & VI-10.1 & Riordino diagrammi dei Casi d'Uso su foglio condiviso LucidSpark & Sofia De Blasi & 2026-04-28 \\[4pt]

\end{longtable}


% --- 6. PROSSIMA RIUNIONE ---
\section{Prossima Riunione}
\begin{itemize}
	\item \textbf{Data:} 2026-04-28
	\item \textbf{Ora:} 14.30
	\item \textbf{OdG preliminare\textsubscript{G}:} 
	\begin{itemize}
		\item Revisione e verifica dei task del terzo sprint;
		\item Allineamento con gli aggiornamenti del referente aziendale;
		\item Avanzamento del Piano di Progetto;
		\item Avanzamento dell'Analisi dei Requisiti;
		\item Pianificazione del quarto sprint e assegnazione dei ruoli.
	\end{itemize} 
\end{itemize}